\chapter{Общий масштаб проекторной суперсвязности}

Рассмотрим максимально благоприятный единый оператор
\begin{equation}
 D_{\rm tot}=D_{\rm space}\widehat\otimes1+gamma\widehat\otimes N/\ell.
\end{equation}
Он объединяет пространственное исчисление и внутреннюю петлю без нового
матричного носителя. Однако тепловое спектральное действие имеет разные
коэффициенты при разных порядках разложения: квадратичный сектор зависит
от $f_2\Lambda^2$, а кривизновый --- от $f_0$. Моменты $f_2$ и $f_0$,
cutoff $\Lambda$ и относительный радиус $\ell$ независимы. Это стандартная
структура heat-kernel expansion и прямое продолжение ранее доказанного
гейта ordinary spectral moment map.

Квадрат суперсвязности не устраняет свободу. При
$\mathbb A=\nabla+a\Phi$ коэффициент $|D\Phi|^2$ меняется как $a^2$, а
квартичные и потенциальные части --- как $a^4$. Нормированный след
фиксирует матричные кратности после выбора $a$, но не выбирает сам $a$.

Следовательно, один product-Dirac делает $E_2$ и $E_4$ совместимыми в
одном действии, но не выводит их относительный размерный коэффициент.
Подстановка планковского масштаба также не решает задачу: проект не вывел,
что внутренний радиус $\ell$ равен планковской длине, а число $1/7$
безразмерно и не может определить этот масштаб.

\begin{theorem}[No-go общего масштаба]
Ни тепловое спектральное действие, ни квадрат текущей суперсвязности не
фиксируют отношение $E_4/E_2$ из данных $N,V_{15}$ и единственного
конечного следа. Конечный радиус остаётся функцией независимого cutoff,
внутренней длины или нормировки нечётного поля.
\end{theorem}

\begin{protocol}
\begin{enumerate}
 \item единый product-Dirac: \textbf{допустим};
 \item независимые моменты $f_2,f_0$: \textbf{остаются};
 \item независимые $\Lambda,\ell$: \textbf{остаются};
 \item свобода $\Phi\mapsto a\Phi$: \textbf{остаётся};
 \item общий масштаб $E_2/E_4$: \textbf{не выведен};
 \item конечный радиус: \textbf{не предсказан};
 \item следующий гейт: \textbf{индуцированное фермионным детерминантом действие}.
\end{enumerate}
\end{protocol}

Машинный сертификат сохранён в
\path{s2t_v5_projector_superconnection_common_scale_gate_results.json}.